\subsection{Negligible sets}

DEFINITION 2. --- Given a measure \(\mu\) on a locally compact space \(X\) , a subset \(A\) of \(X\) is said to be negligible for the measure \(\mu\) if \(|\mu|^*(A) = 0\) .

One also says that A is \(\mu\) -negligible, or simply negligible if no confusion can result. It comes to the same thing to say that the characteristic function \(\varphi_{A}\) is negligible.

PROPOSITION 4. --- Every subset of a negligible set is negligible; every countable union of negligible sets is negligible.

This is an immediate consequence of Props. 1 and 2.

Example. --- Let \(\mu\) be Lebesgue measure on R. Every set \(\{x_{0}\}\) reduced to a point is negligible (cf.~§1, No.~3, Remark 1). It follows that every countable subset of R is negligible for Lebesgue measure. The converse of this proposition is incorrect (§4, Exer. 4 b)).

PROPOSITION 5. --- The complement of the support S of \(\mu\) is the largest negligible open set in X.

For, by Prop. 3, in order that an open set G be negligible, it is necessary and sufficient that \(G \cap S = \varnothing\) , that is, \(G \subset CS\) .

